% =====================================================================
% Related work
% =====================================================================
\section{Related work}
\label{sec:related}

\paragraph{Information bottleneck and conditional compression.}
The information bottleneck (IB) principle formulates representation learning as
retaining information about a target while discarding information about the input
\citep{tishby2000information}. Deep VIB turns this principle into a neural
objective by using a reparameterized stochastic encoder and a variational marginal
prior \citep{alemi2017deep}. Nonlinear and non-parametric bottleneck variants study
alternative variational bounds and distance-based estimates of the information in a
representation \citep{kolchinsky2017nonlinear}. The deterministic-scenario analysis
of \citet{kolchinsky2019caveats} further shows that the usual IB multiplier can fail
to sweep an interior information--accuracy trade-off when labels are deterministic,
motivating the squared-IB functional, which squares the compression term and
restores the interior trade-off.

The conditional entropy bottleneck (CEB) replaces the marginal reference of VIB by a
label-conditional reference and thereby targets conditional residual information
\citep{fischer2020conditional}. Comparisons of variational bounds clarify when this
conditional reference is tighter than a marginal one
\citep{geiger2020variational}. CEB has also been applied to large-scale visual
representations and robustness evaluations \citep{fischer2020robustness,lee2021compressive},
while multivariate extensions retain the same variational bottleneck viewpoint
\citep{abdelaleem2025deep}. These methods leave the relationships among distinct
conditional reference means unspecified. GPB keeps the CEB posterior and its single
conditional KL, but restricts the reference family itself: OPB declares a class
frame and EPB declares a label-distance-preserving axis.

\paragraph{Information bottleneck for robust representations.}
Several works use a bottleneck to suppress nuisance or non-robust information. The
HSIC bottleneck adds layer-wise dependence penalties and analyzes adversarial
robustness \citep{wang2021hsic}. Adversarial Information Bottleneck explicitly
incorporates adversarial objectives \citep{zhai2022adversarial}, and an NLP-specific
variant separates task-relevant robust features from manipulable features
\citep{zhang2022nl_ib}. Causal Information Bottleneck uses a causal decomposition
for the same robust-feature goal \citep{hua2022causal}. These approaches modify the
information criterion, feature selection, or training examples. In contrast, the
geometric constraint in GPB is placed on the conditional prior and is evaluated with
the ordinary CEB KL, so adversarial robustness is improved by the declared geometry
itself, without any additional objective.

\paragraph{Geometric classification representations.}
Neural-collapse studies document an approximately simplex-ETF arrangement of class
means and classifier directions near the terminal phase of training
\citep{papyan2020prevalence,zhu2021geometric_nc}. Other works impose related
geometry directly: orthogonal classifier weights can improve adversarial
robustness \citep{xu2022orthogonal}; ETF classifiers have been used for early exit
\citep{ji2023etf}; and fixed or hierarchy-aware frames have been used to reduce
classifier bias or mistake severity \citep{tong2023fixed_frame}. Prototype-based
constructions date back to equidistant prototype embeddings
\citep{deng2014equidistant}; Neural Collapse by Design trains features and
prototypes jointly on the unit hypersphere so that the global optimum of the NONL
contrastive objective is exactly the simplex ETF \citep{koromilas2026neural},
which we include as a classification-only comparison baseline. The connection
between neural collapse, supervised contrastive learning, and an information
bottleneck has also been analyzed explicitly \citep{wang2023ib_nc}.

These works constrain features, classifier weights, or prototypes through a
classification loss or a fixed readout. OPB addresses a different object: the
label-conditional Gaussian prior in the bottleneck. Its polar/QR frame is recomputed
from the prior encoder, and the anchor-energy classifier reads the same frame without
introducing an independent classifier scale. Thus the geometric class code and the
conditional KL are coupled during training.

\paragraph{Distance-preserving regression and geometric priors.}
Isometric representation learning has traditionally focused on preserving distances
on a data manifold, including deep extensions of isomap-style mappings
\citep{pai2022isometric}, and orthogonal label regression has been used as a
discriminative device in domain adaptation \citep{luo2023doll}; these methods
constrain a deterministic feature map. The regression information-bottleneck
literature, including divergence-based bounds, provides complementary compression
objectives rather than a geometric prior construction \citep{yu2024csd_ib}, and
AdaCap pairs a permutation-based contrastive loss with a Tikhonov closed-form
output layer for small-sample tabular regression \citep{belucci2026adacap}, which
we include as a regression-only comparison baseline. In contrast, EPB declares the
label metric in the mean map of a conditional Gaussian prior: after
standardization the prior means lie on a scaled axis, so latent prior distances
equal the declared label distances, the same axis being the tied regression
readout and the target of the conditional KL. EPB is thus a conditional-prior use
of distance geometry, not a new distance-preserving primitive.

\paragraph{Geometry-aware information bottlenecks.}
Recent work has used the phrase geometric information bottleneck in several distinct
settings. A variational recommender learns a geometric prior for explanation
generation \citep{yan2024geometric_ib}; a data-scarcity formulation penalizes latent
curvature and intrinsic dimension \citep{katende2025vgib}; and activation steering
uses orthogonal subspaces to control language-model interventions
\citep{doan2026geometric_ib}. A concurrent analysis studies the relationship between
information compression and class-wise geometric compression in CEB networks
\citep{adilova2026geometric_compression}. These studies make geometry an important
representation concern, but they do not use the same conditional Gaussian prior for
both compression and prediction, nor do they provide the OPB/EPB alignment-to-
separation construction considered here.

Table~\ref{tab:geometry-boundary} summarizes this object-level distinction: the
geometric primitive is shared with several preceding lines of work, but its role in
the objective is different, as OPB anchors are constrained class centers serving
simultaneously as Gaussian-prior means, KL targets, and readout anchors. OPB is not
an ETF construction: its anchor Gram matrix is $a^2 I_K$ (zero pairwise inner
products), whereas a centered simplex ETF has constant negative off-diagonal inner
products in a $(K-1)$-dimensional subspace.

\begin{table*}[t]
\centering
\footnotesize
\setlength{\tabcolsep}{2pt}
\renewcommand{\arraystretch}{1.12}
\caption{Object-level distinction between the related methods and GPB. IB KL:
the method's constraint enters an information-bottleneck KL term; Tied readout:
the readout is tied to the constrained geometry; Fixed / hierarchy: fixed or
hierarchy-aware frames; Isometric maps: isometric feature maps.}
\label{tab:geometry-boundary}
\begin{tabular}{>{\centering\arraybackslash}m{2.2cm} >{\centering\arraybackslash}m{2.4cm} >{\centering\arraybackslash}m{1.0cm} >{\centering\arraybackslash}m{1.6cm} >{\centering\arraybackslash}m{2.8cm} p{2.9cm}}
\hline
Approach & Declared geometry & IB KL & Tied readout & Stochastic posterior & Constrained object \\
\hline
Prototype methods & $\times$ & $\times$ & $\times$ (partial) & $\times$ (mostly) & prototypes or features \\
ETF / orthogonal & \checkmark & $\times$ & \checkmark & $\times$ (mostly) & classifier weights or feature directions \\
Fixed / hierarchy & \checkmark & $\times$ & varies & $\times$ (mostly) & feature coordinates or a frame \\
Isometric maps & \checkmark & $\times$ & varies & $\times$ (mostly) & feature map \\
CEB & $\times$ & \checkmark & $\times$ & \checkmark & label-conditional reference \\
GPB (OPB / EPB) & \checkmark & \checkmark & \checkmark & \checkmark & label-conditional Gaussian prior means \\
\hline
\end{tabular}
\end{table*}
